\documentclass[11pt]{article}
\usepackage[a4paper,margin=24mm]{geometry}
\usepackage[T1]{fontenc}
\usepackage{lmodern,amsmath,amssymb,amsthm,booktabs,microtype}
\usepackage[colorlinks=true,urlcolor=blue,linkcolor=blue,citecolor=blue]{hyperref}
\newcommand{\tr}{\operatorname{tr}}
\newcommand{\sech}{\operatorname{sech}}
\newtheorem{proposition}{Proposition}
\title{Selected preparations and the Rindler time kernel\\
\large Exact source response and a semiclassical shape correction}
\author{Alan Day \and Chengheng\thanks{AI research assistant (OpenAI).}}
\date{8 November 2026\quad Local research note}
\begin{document}\maketitle
\begin{abstract}
The fixed-charge bilocal construction gives access to observables
between controlled selected preparations. We express the single-selected
functional as the second response of a normalized source state. On the
charged subalgebra the source is the physical Hamiltonian. Applying the
known thermal chord saddle predicts a first-order change in the
curvature two-point function which cannot be absorbed into a constant
rescaling of time. Exact gap-moment expansions through time degree
fourteen and finite-deformation calculations test this prediction.
The source identity is exact; differentiating the semiclassical saddle
requires an additional uniform remainder estimate not proved here.
\end{abstract}

\section{Physical question, normalization and existing input}
The Susskind program asks for a quantitative comparison of derivative
singlets with the bulk meson spectrum \cite{mss}. Its singlets are
$\sum_i\psi_i^\dagger\partial_t^n\psi_i$. Our earlier work has
identified their nonconserved components and reconstructed the actual
centered bilocal in a fixed charge sector \cite{bilocal}. Before using
its time spectrum in a bulk dictionary, we ask how the choice of
preparation changes that spectrum.

Use the standard overlapping complex Gaussian p-body CAR model,
the same disorder in all fields, $\hbar=1$, string time t, and
$p^2/N\to\lambda>0$ at half filling. First take this double-scaled
limit. Subsequent small-$\lambda$ statements keep string time fixed.
Write
\[
 \varepsilon=4\lambda,\quad q=e^{-\varepsilon},\quad
 C=J_s^2/\lambda,\quad e=H/\sqrt C,\quad
 K=H/(2J_s)=e/\sqrt\varepsilon,\quad u=2J_st.
\]
The energy e has unit variance and K has variance $1/\varepsilon$.
Let x be the normalized nonconserved curvature singlet of
\cite{selected}. In its limiting tracial algebra,
$\tau(x^2)=1$ and its two-energy kernel is the q-Mehler kernel
with parameter $r=q^2$. Define
\[
 x(u)=e^{iuK}xe^{-iuK},\quad
 C_0(u)=\tau(x(u)x),\qquad C_1(u)=\sigma_1(x(u)x).
\]
The previously proved single-selected restriction is
\begin{equation}\label{density}
 \sigma_1(Y)=\tau(f_q(e)Y),\qquad
 f_q(e)=1-\sqrt q+\sqrt q\,e^2,
 \quad Y\in W^*(e,x).
\end{equation}
It is normalized, positive and stationary on this subalgebra.
The corresponding finite-N statement is the sequential limit of
actual exactly pinched observables in the specified preparation.
No rate uniform in $\lambda\to0$ is assumed.

\newpage
\section{Exact source-response identities}
For a centered self-adjoint A in a tracial algebra define
\[
 F_A(z;Y)=\frac{\tau(e^{-zA}Y)}{\tau(e^{-zA})}.
\]
For bounded A this is analytic near zero. Direct differentiation,
with $\tau(A)=0$, gives
\begin{equation}\label{source}
 F_A''(0;Y)=\tau(A^2Y)-\tau(A^2)\tau(Y).
\end{equation}
No commutation of A and Y is needed. This is the usual normalized
source susceptibility, applied here to the known selected density.

For each fixed $q<1$, K is bounded. Let
\begin{equation}\label{thermal}
 G_{\varepsilon,\beta}(u)=
 \frac{\tau(e^{-\beta K}x(u)x)}{Z_\varepsilon(\beta)},
 \qquad Z_\varepsilon(\beta)=\tau(e^{-\beta K}).
\end{equation}
Here $\beta$ is dimensionless and conjugate to K; physical inverse
string energy is $\beta/(2J_s)$. This is a limiting tracial Gibbs
functional, not an assumed bulk Rindler state or an assertion about
quenched finite-N normalization.
\begin{proposition}
At every fixed $\varepsilon>0$ and real u,
\begin{equation}\label{exact}
 C_1(u)-C_0(u)=e^{-\varepsilon/2}\varepsilon
   \left.\partial_\beta^2G_{\varepsilon,\beta}(u)\right|_{\beta=0}.
\end{equation}
\end{proposition}
\begin{proof}
Equation \eqref{density} gives the difference as
$\sqrt q[\varepsilon\tau(K^2x(u)x)-\tau(x(u)x)]$.
Apply \eqref{source} and $\tau(K^2)=1/\varepsilon$.
\end{proof}

The full covariant tower contains a size field W and has the
different density established in \cite{selected,bilocal}:
\[
 \rho_1=\sqrt q\,(e-W)^2+[1-\sqrt q(1+\lambda)]I,
 \quad\tau((e-W)^2)=1+\lambda.
\]
Consequently, for the controlled integrable covariant observables,
\begin{equation}\label{full}
 \sigma_1(Y)-\tau(Y)=\sqrt q\,F_{e-W}''(0;Y).
\end{equation}
At fixed q, e and W are bounded. The actual bilocal limit has all
fixed moments, so the differentiation is also valid for those
integrable Y. This identity retains the terms
$-\tau((eW+We)Y)+\tau((W^2-\lambda)Y)$ absent from the charged
restriction. The source $e-W$ is not the physical time generator.
It supplies a source response for the full bilocal, not a thermal
interpretation of its generally nonstationary selected state.

\newpage
\section{The thermal saddle and a testable shape correction}
The leading thermal matter correlator is known from large-p SYK
and the semiclassical chord theory \cite{ls}. In its conventions,
put $\theta=\pi v/2$; the thermal equation and real-time continuation are
\begin{equation}\label{saddle}
 \beta=\frac{2\theta}{\cos\theta},\qquad
 g_{\beta,d}(u)=
 \left[\frac{\cos\theta}{\cosh(u\cos\theta+i\theta)}\right]^{2d}.
\end{equation}
The Hamiltonian in these units is K, whose variance is
$1/\varepsilon$. Our curvature channel has $d=2$ because its
Hamiltonian crossing weight is $q^2$, not because its bulk mass
has been identified. Equation \eqref{saddle} is the unsplit
thermal correlator; a midpoint thermal correlator has different
source derivatives.

Since $\theta=\beta/2+O(\beta^3)$, with $T=\tanh u$,
\[
 \log g_{\beta,d}
 =-2d\log\cosh u-id\beta T
       +\frac{d\beta^2}{4}(uT-T^2)+O(\beta^3).
\]
Therefore
\begin{equation}\label{response}
 \left.\partial_\beta^2g_{\beta,d}\right|_0
 =\frac d2\sech^{2d}u\,[u\tanh u-(1+2d)\tanh^2u].
\end{equation}
Combining \eqref{response} with the exact \eqref{exact} predicts
\begin{equation}\label{prediction}
 \frac{C_1(u)-C_0(u)}{\varepsilon}\longrightarrow
 R(u):=\sech^4u\,[u\tanh u-5\tanh^2u].
\end{equation}
The precise sufficient hypothesis is convergence of
$\partial_\beta^2G_{\varepsilon,\beta}|_0$ to
$\partial_\beta^2g_{\beta,2}|_0$, locally uniformly in u.
A locally uniform analytic saddle expansion in a complex
$\beta$ neighborhood would imply it by Cauchy's formula.
Pointwise thermal convergence alone does not. We do not prove
this differentiated remainder bound in this note, and label
\eqref{prediction} as a semiclassical prediction rather than a
new all-time convergence theorem.

Nevertheless the short-time coefficient has an independent exact
check. Mehler conditional moments give
\begin{equation}\label{curvature}
 -(C_1-C_0)''(0)=
 \frac{e^{-\varepsilon/2}(1+e^{-\varepsilon})
                   (1-e^{-2\varepsilon})^2}{\varepsilon}
 =8\varepsilon+O(\varepsilon^2).
\end{equation}
The prediction has $R(u)=-4u^2+11u^4+O(u^6)$.
It cannot be a constant time recalibration:
$a u(\sech^4u)'=-4a u^2+(28a/3)u^4+O(u^6)$.
Matching the quadratic term fixes $a=1$, which fails at the
quartic term. An overall amplitude is fixed by $C_1(0)=C_0(0)=1$.

\newpage
\section{Independent coefficients and finite-deformation evidence}
Let $Q_j$ be the monic unit-variance q-Hermite polynomials,
$e^a=\sum_jb_{aj}Q_j(e)$. The exact two-energy moments are
\[
 T_{ab}(q)=\sum_jb_{aj}b_{bj}[j]_q!\,q^{2j}.
\]
Thus the selected-minus-trace gap moment before physical time
rescaling is
\[
 D_n(q)=\sum_{a=0}^n(-1)^{n-a}\binom na
                 [T_{a+2,n-a}(q)-T_{a,n-a}(q)].
\]
All quantities here are integer polynomials in q. If
$D_{2k}(q)=\sum_jd_jq^j$, the coefficient of $\varepsilon^\ell$
in $e^{-\varepsilon/2}D_{2k}(e^{-\varepsilon})$ is exactly
\[
 \frac1{\ell!}\sum_jd_j(-j-1/2)^\ell.
\]
The physical time coefficient includes
$(-1)^k/[(2k)!\varepsilon^k]$.
For $2k=0,2,\ldots,14$, all 36 terms with $\ell\le k$ vanish
exactly. The first surviving coefficients reproduce
\[
 R(u)=-4u^2+11u^4-\frac{769}{45}u^6+
 \frac{6236}{315}u^8-\frac{1201}{63}u^{10}
 +\frac{7551623}{467775}u^{12}
 -\frac{529250077}{42567525}u^{14}+O(u^{16}).
\]
These are finite exact checks, not a proof for all Taylor orders.

For real time use the Jacobi matrix T with off-diagonal
$\sqrt{[n]_q}$ and $D=\operatorname{diag}(q^{2n})$.
With $a=e^{iuT/\sqrt\varepsilon}|0\rangle$ and
$b=e^{iuT/\sqrt\varepsilon}|2\rangle$, the exact infinite-matrix
formulas inherited from \cite{selected} are
\[
 C_0=\langle a,Da\rangle,\qquad
 C_1-C_0=\sqrt{q(1+q)}\langle b,Da\rangle.
\]
The difference is computed directly. On the seven-point grid
$u=0,0.25,\ldots,1.5$ the finite-deformation results are
\begin{center}\begin{tabular}{rrr}\toprule
$\varepsilon$ & Jacobi dimensions & $\max| (C_1-C_0)/\varepsilon-R|$\\\midrule
0.08 & 214,278 & 0.077469\\
0.04 & 364,428 & 0.040022\\
0.02 & 664,728 & 0.020335\\
0.01 & 1264,1328 & 0.010249\\\bottomrule
\end{tabular}\end{center}
The first dimension uses $\lceil12/\varepsilon\rceil+64$.
The displayed differences measure finite-deformation bias, not
numerical error. The parameter $\lambda$ equals $\varepsilon/4$.

\newpage
\section{Error accounting and the next construction}
For $z=2|u|/\sqrt{\varepsilon(1-q)}$, a common exact-arithmetic
packet truncation bound inherited from \cite{selected} is
\[
 4[1+\sqrt{q(1+q)}]e^z\frac{z^{L-2}}{(L-2)!}.
\]
The saved parameters satisfy $z^2\le L(L-1)$; all nonzero-time
bounds are below $10^{-25}$. This bounds truncation, not floating
roundoff. We compare L and L+64, retain their difference after
division by $\varepsilon$, check packet normalization, and at
$\varepsilon=0.08$ compare sparse exponential action against an
independent symmetric tridiagonal eigensolver. All numerical
comparisons use float64 and empirical tolerances, not intervals.

The largest summand-cancellation ratio is about 414. For example,
at $\varepsilon=0.01,u=0.25$ the coefficient differs from the
prediction by about $2.68\%$, despite negligible cutoff change.
At $u=0.5$ its error is about $1.98\%$. These differences must not
be relabeled as numerical uncertainty or omitted by an extrapolation.

The physical outcome is a source-response method for handling
preparation dependence while retaining the actual microscopic
observable. Under the differentiated saddle hypothesis, the
trace and selected curvature have the same leading Rindler
kernel, but a distinct first correction. This is state response,
not a change in an invariant particle mass.

The next analytic step is to control the twice-differentiated
thermal q-Hermite integral near $\beta=0$, with fixed compact u.
The full bilocal additionally needs the mixed e,W source in
\eqref{full}; using \eqref{exact} there would drop real neutral
terms. Spatial translations, mass generators, arbitrary selected
interlacings and finite-G dynamics remain separate unfinished
parts of the intended dictionary. This stage introduces no new
Git publication or issue reply.

\begin{thebibliography}{9}\small
\bibitem{mss} S. Miyashita, Y. Sekino and L. Susskind,
\emph{Holograms and Standard Models},
\href{https://arxiv.org/html/2607.05678v2}{arXiv:2607.05678v2},
Secs.~4.1--4.3. The mass-spectrum comparison is an open quantitative
test; a boost kernel alone does not determine that spectrum.
\bibitem{ls} H. W. Lin and D. Stanford, \emph{A symmetry algebra
in double-scaled SYK},
\href{https://arxiv.org/html/2307.15725}{arXiv:2307.15725},
Eqs.~(137)--(138), Sec.~4.3 and Appendix H, Eqs.~(409)--(410).
Their chord deformation parameter is our $\varepsilon$.
\bibitem{selected} A. Day and Chengheng, \emph{Single-selected
covariant density and curvature dynamics}, 8 November 2026;
\texttt{manuscript/susskind-slater-selected-state/main.tex}.
\bibitem{bilocal} A. Day and Chengheng, \emph{Fixed-charge bilocals
and selected preparations}, 8 November 2026;
\texttt{manuscript/susskind-sector-bilocal/main.tex}.
\bibitem{target} A. Day and Chengheng, bulk meson target and
translation notes, 6 November 2026;
\texttt{notes/susskind-bulk-target.md},
\texttt{notes/susskind-translations.md}.
\end{thebibliography}
\end{document}
